\chapter{คลังรหัสและองค์ประกอบสกิลวิศวกรรมข้อมูลคลาวด์และแมชชีนเลิร์นนิง}

ภาคผนวกนี้รวบรวมข้อมูลจำเพาะ โค้ดต้นแบบ และสถาปัตยกรรมระดับองค์กรของสกิล BigQuery, Dataform, dbt, Spark Serverless และการสร้างโมเดลการเรียนรู้ของเครื่องประจำเล่มที่ 5

\section{สกิล: bigquery-sql}

\begin{teachingplan}{ข้อมูลจำเพาะและเมทาดาทาของสกิล}
\textbf{ชื่อสกิล:} \texttt{bigquery-sql}\\
\textbf{วัตถุประสงค์:} ออกแบบและปรับแต่งคิวรี SQL บน BigQuery เพื่อลดค่าใช้จ่ายและเร่งความเร็วประมวลผล\\
\textbf{เทคนิคสำคัญ:} Partition Pruning, Cluster Filtering, Window Functions over Self-Joins
\end{teachingplan}

\subsection{โค้ดแกนกลาง: แม่แบบการสร้างตาราง Partitioned และ Clustered}

\begin{lstlisting}[language=SQL, caption={BigQuery DDL Optimized Table Creation}]
CREATE OR REPLACE TABLE `my_project.soil_analytics.sensor_telemetry`
(
  device_id STRING NOT NULL,
  timestamp TIMESTAMP NOT NULL,
  ph_value FLOAT64,
  moisture_percent FLOAT64,
  temperature_celsius FLOAT64,
  battery_level FLOAT64
)
PARTITION BY DATE(timestamp)
CLUSTER BY device_id
OPTIONS(
  description = "Partitioned daily & clustered by device_id to minimize slot scan costs"
);
\end{lstlisting}

\section{สกิล: dataform-bigquery}

\begin{teachingplan}{ข้อมูลจำเพาะและเมทาดาทาของสกิล}
\textbf{ชื่อสกิล:} \texttt{dataform-bigquery}\\
\textbf{วัตถุประสงค์:} พัฒนา Data Transformation Pipeline ด้วย Dataform SQLX พร้อม Assertion Tests\\
\textbf{สถาปัตยกรรม:} Medallion Layers (Bronze $\rightarrow$ Silver $\rightarrow$ Gold)
\end{teachingplan}

\subsection{โค้ดแกนกลาง: โมเดล Dataform SQLX พร้อมการตรวจสอบคุณภาพ}

\begin{lstlisting}[language=SQL, caption={Dataform SQLX Table with Assertion Constraints}]
config {
  type: "table",
  schema: "silver_layer",
  description: "Cleaned and calibrated soil measurement data",
  assertions: {
    nonNull: ["device_id", "timestamp"],
    rowConditions: [
      "ph_value >= 0.0 AND ph_value <= 14.0",
      "moisture_percent >= 0.0 AND moisture_percent <= 100.0"
    ]
  }
}

SELECT
  device_id,
  timestamp,
  ROUND(ph_value, 2) AS ph_value,
  ROUND(moisture_percent, 1) AS moisture_percent,
  temperature_celsius
FROM
  ${ref("raw_sensor_telemetry")}
WHERE
  ph_value IS NOT NULL
\end{lstlisting}

\section{สกิล: gcp-spark}

\begin{teachingplan}{ข้อมูลจำเพาะและเมทาดาทาของสกิล}
\textbf{ชื่อสกิล:} \texttt{gcp-spark}\\
\textbf{วัตถุประสงค์:} พัฒนาโค้ด PySpark บน Dataproc Serverless สำหรับประมวลผล Big Data และ ML\\
\textbf{เทคนิคสำคัญ:} Broadcast Joins, Off-Heap Memory Configuration, Spark ML Pipeline
\end{teachingplan}

\subsection{โค้ดแกนกลาง: PySpark MLlib Pipeline บน Dataproc Serverless}

\begin{lstlisting}[language=Python, caption={PySpark Random Forest Regression Pipeline}]
from pyspark.sql import SparkSession
from pyspark.ml.feature import VectorAssembler
from pyspark.ml.regression import RandomForestRegressor
from pyspark.ml import Pipeline

spark = SparkSession.builder.appName("SoilMoisturePrediction").getOrCreate()

# Load cleaned telemetry from BigQuery connector
df = spark.read.format("bigquery") \
    .option("table", "my_project.silver_layer.sensor_telemetry") \
    .load()

feature_cols = ["ph_value", "temperature_celsius", "battery_level"]
assembler = VectorAssembler(inputCols=feature_cols, outputCol="features")

rf = RandomForestRegressor(featuresCol="features", labelCol="moisture_percent", numTrees=50)
pipeline = Pipeline(stages=[assembler, rf])

model = pipeline.fit(df)
spark.stop()
\end{lstlisting}


\section{ดัชนีคลังไฟล์ข้อกำหนด SKILL.md ประจำเล่มที่ 5}

หนังสือเล่มนี้ได้จัดเก็บไฟล์ข้อกำหนด \texttt{SKILL.md} และองค์ประกอบทั้งหมดไว้ในโฟลเดอร์ \texttt{skills/} ซึ่งผู้อ่านสามารถนำไปเรียกใช้งานกับ AI Agent หรือศึกษาโครงสร้างเชิงลึกได้ทันที:

\small
\begin{longtable}{|c|p{0.45\textwidth}|p{0.40\textwidth}|}
\hline
\textbf{ลำดับ} & \textbf{ชื่อสกิล (Skill Name)} & \textbf{ที่อยู่ไฟล์ต้นฉบับ} \\ \hline
\endfirsthead
\hline
\textbf{ลำดับ} & \textbf{ชื่อสกิล (Skill Name)} & \textbf{ที่อยู่ไฟล์ต้นฉบับ} \\ \hline
\endhead
1 & \texttt{bigquery-data-transfer-service} & \texttt{skills/bigquery-data-transfer-service/SKILL.md} \\ \hline
2 & \texttt{bigquery-sql} & \texttt{skills/bigquery-sql/SKILL.md} \\ \hline
3 & \texttt{data-autocleaning} & \texttt{skills/data-autocleaning/SKILL.md} \\ \hline
4 & \texttt{dataform-bigquery} & \texttt{skills/dataform-bigquery/SKILL.md} \\ \hline
5 & \texttt{dbt-bigquery} & \texttt{skills/dbt-bigquery/SKILL.md} \\ \hline
6 & \texttt{discovering-gcp-data-assets} & \texttt{skills/discovering-gcp-data-assets/SKILL.md} \\ \hline
7 & \texttt{enforcing-resource-attribution} & \texttt{skills/enforcing-resource-attribution/SKILL.md} \\ \hline
8 & \texttt{gcp-data-pipelines} & \texttt{skills/gcp-data-pipelines/SKILL.md} \\ \hline
9 & \texttt{gcp-dataflow} & \texttt{skills/gcp-dataflow/SKILL.md} \\ \hline
10 & \texttt{gcp-pipeline-orchestration} & \texttt{skills/gcp-pipeline-orchestration/SKILL.md} \\ \hline
11 & \texttt{gcp-pipeline-resource-provisioning} & \texttt{skills/gcp-pipeline-resource-provisioning/SKILL.md} \\ \hline
12 & \texttt{gcp-spark} & \texttt{skills/gcp-spark/SKILL.md} \\ \hline
13 & \texttt{gcp-spark-troubleshooting} & \texttt{skills/gcp-spark-troubleshooting/SKILL.md} \\ \hline
14 & \texttt{google-cloud-auth-verification} & \texttt{skills/google-cloud-auth-verification/SKILL.md} \\ \hline
15 & \texttt{managing-python-dependencies} & \texttt{skills/managing-python-dependencies/SKILL.md} \\ \hline
16 & \texttt{ml-best-practices} & \texttt{skills/ml-best-practices/SKILL.md} \\ \hline
17 & \texttt{notebook-guidance} & \texttt{skills/notebook-guidance/SKILL.md} \\ \hline
18 & \texttt{schema-mapping} & \texttt{skills/schema-mapping/SKILL.md} \\ \hline
19 & \texttt{uv} & \texttt{skills/uv/SKILL.md} \\ \hline
\end{longtable}
\normalsize
